%---------------------------------------------------------------------
% Part one: the Makki period. Five questions.
%---------------------------------------------------------------------

\section*{حصہ اول --- مکی دور}
\markboth{مکی دور}{}

%---------------------------------------------------------------------
\begin{sawal}{سوال \textenglish{1} \ \textnormal{\small(چاروں پرچوں میں --- بہار \textenglish{2013} س\,\textenglish{2}، خزاں \textenglish{2012} س\,\textenglish{2}، خزاں \textenglish{2017} س\,\textenglish{1}(الف\textenglish{2})، بہار \textenglish{2022} س\,\textenglish{1})}}
سیرت طیبہ کا لغوی و اصطلاحی مفہوم، مطالعۂ سیرت کی ضرورت و اہمیت، اور سیرت
طیبہ کے مآخذ و مصادر تفصیل سے بیان کریں۔
\end{sawal}

\begin{hal}
یہ سوال \textbf{ہر پرچے میں} کسی نہ کسی صورت میں آیا ہے --- سب سے زیادہ
دہرایا جانے والا موضوع۔

\medskip
\textbf{لغوی و اصطلاحی مفہوم۔} ''سیرت'' عربی لفظ ہے، معنی: طرزِ عمل، چال ڈھال۔
اصطلاح میں: نبی کریم \saw کی حیاتِ طیبہ کے حالات، اقوال، افعال اور تقریرات
(جن باتوں کی آپ \saw نے تصدیق کی) کا مجموعہ۔

\medskip
\textbf{ضرورت و اہمیت۔} نبی کریم \saw \textbf{اسوۂ حسنہ} (بہترین نمونہ) ہیں:
\ar{لَقَدْ كَانَ لَكُمْ فِي رَسُولِ اللَّهِ أُسْوَةٌ حَسَنَةٌ} (الاحزاب،
\textenglish{21})۔ سیرت کا مطالعہ تین وجوہ سے ضروری ہے: \textbf{(۱)} قرآن کی
عملی تفسیر --- قرآن کے احکام کا زندہ نمونہ سیرت میں ملتا ہے؛ \textbf{(۲)}
اطاعتِ رسول ایمان کی شرط ہے؛ \textbf{(۳)} انفرادی و اجتماعی زندگی کے ہر
شعبے --- عبادت، معیشت، سیاست، خاندان --- میں براہِ راست رہنمائی ملتی ہے۔

\medskip
\textbf{مآخذ و مصادر --- چار درجے:}
\begin{enumerate}\setlength{\itemsep}{2pt}
  \item \textbf{قرآن مجید} --- بنیادی اور یقینی ترین مآخذ؛ سیرت کے متعدد
        واقعات (بدر، احد، ہجرت، اسریٰ) کا براہِ راست ذکر۔
  \item \textbf{حدیث و سنت} --- نبی \saw کے اقوال و افعال کی مستند روایات،
        خصوصاً صحاح ستہ۔
  \item \textbf{کتبِ سیرت} --- ابتدائی اور مستند: \textbf{سیرتِ ابن اسحاق}
        (م \textenglish{151}ھ، سب سے قدیم مربوط سیرت، بعد میں ابن ہشام نے
        مرتب کی)، \textbf{سیرتِ ابن ہشام}، \textbf{طبقاتِ ابن سعد}،
        \textbf{تاریخِ طبری}۔
  \item \textbf{اشعارِ عرب اور غزوات کی روایات (مغازی)} --- زمانۂ نزول کے
        حالات سمجھنے میں مددگار، جیسے واقدی کی \textbf{کتاب المغازی}۔
\end{enumerate}
\end{hal}

%---------------------------------------------------------------------
\begin{sawal}{سوال \textenglish{2} \ \textnormal{\small(خزاں \textenglish{2017}، سوال \textenglish{2})}}
حضور \saw کی ولادت سے قبل عرب کے حالات پر جامع نوٹ تحریر کریں۔
\end{sawal}

\begin{hal}
اس دور کو \textbf{زمانۂ جاہلیت} کہا جاتا ہے --- علم کی کمی کے معنی میں نہیں،
بلکہ اخلاقی و عقائدی گمراہی کے معنی میں۔

\medskip
\textbf{عقائد۔} بت پرستی عام تھی --- خانہ کعبہ میں \textenglish{360} بت
نصب تھے۔ چند قبائل یہودیت، عیسائیت یا حنیفیت (ابراہیمی توحید) پر بھی قائم
تھے۔

\medskip
\textbf{معاشرتی حالات۔} \textbf{قبائلی عصبیت} فیصلہ کن قوت تھی --- معمولی
باتوں پر برسوں جاری رہنے والی جنگیں (جنگِ بسوس، جنگِ داحس و غبراء)۔
\textbf{عورت کا مقام} پست تھا --- بیٹیوں کو زندہ درگور کیا جاتا
(\ar{وَإِذَا الْمَوْءُودَةُ سُئِلَتْ}، التکویر، \textenglish{8})۔ سود،
شراب نوشی اور جوا عام تھے۔

\medskip
\textbf{مثبت پہلو --- بالکل تاریک نہیں تھا۔} مہمان نوازی، وعدہ نبھانا، اور
شعر و ادب میں کمال ان کی نمایاں خوبیاں تھیں --- مکہ کا سالانہ میلہ
\textbf{سوقِ عکاظ} شاعری کے مقابلوں کا مرکز تھا۔ \textbf{حلفِ فضول} جیسا
معاہدہ (مظلوموں کی مدد کا عہد) بھی اسی دور کی پیداوار ہے، جس میں نبی \saw
نے جوانی میں شرکت کی اور بعد میں نبوت کے بعد بھی اس کی تعریف فرمائی۔

\medskip
\textbf{خلاصہ۔} عرب معاشرہ اخلاقی طور پر زوال پذیر مگر بالکل بے ترتیب نہیں
تھا --- بعثتِ نبوی نے اسی خام مادے کو ایک عالمگیر امت میں بدل دیا۔
\end{hal}

%---------------------------------------------------------------------
\begin{sawal}{سوال \textenglish{3} \ \textnormal{\small(بہار \textenglish{2013}، سوال \textenglish{3} اور بہار \textenglish{2022}، سوال \textenglish{3})}}
بعثتِ نبوی اور دعوتِ اسلام کے آغاز پر تفصیلی نوٹ لکھیں، نیز دعوت و تبلیغ کے
اہم اصول قلمبند کریں۔
\end{sawal}

\begin{hal}
\textbf{بعثت۔} چالیس سال کی عمر میں، \textenglish{610}ء میں، غارِ حرا میں
پہلی وحی نازل ہوئی: \ar{اقْرَأْ بِاسْمِ رَبِّكَ الَّذِي خَلَقَ} (العلق،
\textenglish{1})۔ اس کے بعد \textbf{فترتِ وحی} (وحی کا کچھ عرصہ وقفہ) آئی،
پھر سورۃ المدثر سے سلسلہ دوبارہ شروع ہوا: \ar{يَا أَيُّهَا الْمُدَّثِّرُ
قُمْ فَأَنذِرْ} --- اے کملی اوڑھنے والے! اٹھ اور ڈرا۔

\medskip
\textbf{دعوت کے تین مراحل:}
\begin{enumerate}\setlength{\itemsep}{2pt}
  \item \textbf{خفیہ دعوت} (تین سال) --- قریبی رشتہ داروں اور دوستوں تک؛
        سب سے پہلے حضرت خدیجہؓ، حضرت علیؓ، حضرت ابوبکرؓ اور حضرت زیدؓ ایمان
        لائے۔ تعلیم و تربیت کا مرکز \textbf{دار ارقم} (صفا کے قریب حضرت
        ارقمؓ کا گھر) بنا۔
  \item \textbf{علانیہ دعوت --- خاص الرشتہ داروں میں} --- سورۃ الشعراء
        \textenglish{214} کے نزول پر: \ar{وَأَنذِرْ عَشِيرَتَكَ
        الْأَقْرَبِينَ}، کوہِ صفا پر خطاب۔
  \item \textbf{علانیہ دعوت --- عام لوگوں میں} --- سورۃ الحجر
        \textenglish{94} کے حکم پر: \ar{فَاصْدَعْ بِمَا تُؤْمَرُ}۔
\end{enumerate}

\medskip
\textbf{دعوت و تبلیغ کے اہم اصول (قرآن سے ماخوذ):}
\begin{itemize}\setlength{\itemsep}{1pt}
  \item \textbf{حکمت اور اچھی نصیحت:} \ar{ادْعُ إِلَىٰ سَبِيلِ رَبِّكَ
        بِالْحِكْمَةِ وَالْمَوْعِظَةِ الْحَسَنَةِ} (النحل، \textenglish{125})۔
  \item \textbf{نرمی اور بہترین طریقے سے مباحثہ:} \ar{وَجَادِلْهُم بِالَّتِي
        هِيَ أَحْسَنُ} (وہی آیت)۔
  \item \textbf{صبر و استقامت} --- تیرہ سالہ مکی دور کی مسلسل مخالفت کے
        باوجود دعوت جاری رہی۔
  \item \textbf{عمل سے پہلے خود عمل} --- داعی کا اپنا کردار سب سے بڑا وسیلہ۔
\end{itemize}
\end{hal}

%---------------------------------------------------------------------
\begin{sawal}{سوال \textenglish{4} \ \textnormal{\small(خزاں \textenglish{2017}، سوال \textenglish{3})}}
حضرت حمزہؓ اور حضرت عمرؓ کے قبولِ اسلام کے واقعات تحریر کریں۔
\end{sawal}

\begin{hal}
دونوں واقعات \textbf{بعثت کے چھٹے سال} پیش آئے اور مسلمانوں کے لیے بڑی طاقت
کا باعث بنے --- حدیث میں دعا ہے: ''اے اللہ! ان دو عمر (ابو جہل یا عمر بن
خطاب) میں سے جو تجھے زیادہ محبوب ہو، اس سے اسلام کو قوت دے'' --- اور قبولیت
حضرت عمرؓ کے حق میں ہوئی۔

\medskip
\textbf{حضرت حمزہؓ۔} ابو جہل نے صفا کے قریب نبی کریم \saw کو ایذا دی۔ آپ
\saw کے چچا حضرت حمزہؓ، جو اس وقت مسلمان نہ تھے، شکار سے واپسی پر یہ سن کر
غصے میں آئے اور کعبہ جا کر ابو جہل کے سر پر کمان دے ماری، پھر اعلان کیا:
''میں بھی محمد \saw کے دین پر ہوں، جو کچھ وہ کہتے ہیں میں اسے مانتا ہوں''۔
یہ اعلانِ اسلام تھا، محض غصے کا ردِعمل نہیں --- اسی لمحے سے حمزہؓ کا ایمان
پختہ ہو گیا اور وہ ''اسد اللہ'' (اللہ کا شیر) کے لقب سے مشہور ہوئے۔

\medskip
\textbf{حضرت عمرؓ۔} پہلے سخت مخالف تھے، نبی کریم \saw کو قتل کرنے کے ارادے
سے نکلے۔ راستے میں معلوم ہوا کہ ان کی اپنی بہن حضرت فاطمہؓ اور بہنوئی سعید
بن زیدؓ مسلمان ہو چکے ہیں۔ غصے میں ان کے گھر پہنچے، بہن کو مارا، پھر جس
کاغذ پر سورۃ طٰہٰ کی آیات لکھی تھیں وہ پڑھیں --- قرآن کی تاثیر نے دل بدل
دیا، اور وہ سیدھے دار ارقم پہنچ کر نبی کریم \saw کے سامنے اسلام قبول کر
لیا۔ اسی دن مسلمانوں نے پہلی بار \textbf{علانیہ} خانہ کعبہ میں نماز ادا کی۔

\medskip
\textbf{اہمیت۔} دونوں شخصیات کے اسلام لانے سے مسلمانوں کو حوصلہ اور عزت
ملی --- حضرت عبد اللہ بن مسعودؓ کے قول کے مطابق، حضرت عمرؓ کے اسلام لانے
تک مسلمان خفیہ رہتے تھے، ان کے بعد علانیہ عبادت کرنے لگے۔
\end{hal}

%---------------------------------------------------------------------
\begin{sawal}{سوال \textenglish{5} \ \textnormal{\small(بہار \textenglish{2013}، سوال \textenglish{4})}}
ہجرتِ حبشہ پر نوٹ لکھیں، نیز مکی دور میں مسلمانوں کو درپیش مصائب کا جائزہ
پیش کریں۔
\end{sawal}

\begin{hal}
\textbf{مصائب کی نوعیت۔} جیسے جیسے دعوت پھیلی، قریش کی مخالفت شدت اختیار
کرتی گئی --- کمزور و غلام صحابہ پر تشدد (حضرت بلالؓ، حضرت عمارؓ اور ان کے
والدین حضرت یاسرؓ و سمیہؓ --- سمیہؓ اسلام کی پہلی شہید)، معاشرتی و معاشی
مقاطعہ، اور خود نبی کریم \saw کی توہین و اذیت (جیسے راستے میں کانٹے بچھانا،
سجدے میں اوجھڑی ڈالنا)۔

\medskip
\textbf{ہجرتِ حبشہ --- دو مرحلے۔} جب ظلم ناقابلِ برداشت ہو گیا تو نبی کریم
\saw نے حکم دیا کہ کمزور مسلمان \textbf{حبشہ} (موجودہ ایتھوپیا) ہجرت کر
جائیں، جہاں کا عیسائی بادشاہ \textbf{نجاشی} عادل مشہور تھا۔

\begin{itemize}\setlength{\itemsep}{1pt}
  \item \textbf{پہلی ہجرت} (\textenglish{5} نبوی) --- گیارہ مرد اور چار
        عورتیں، قیادت حضرت عثمانؓ بن عفان اور ان کی اہلیہ حضرت رقیہؓ (نبی
        \saw کی صاحبزادی)۔
  \item \textbf{دوسری ہجرت} (\textenglish{6} نبوی) --- تعداد بڑھ کر تقریباً
        \textenglish{83} مرد و عورت، قیادت حضرت جعفر بن ابی طالبؓ۔
\end{itemize}

\medskip
\textbf{نجاشی کے دربار میں مکالمہ۔} قریش نے وفد بھیج کر مہاجرین کی واپسی کا
مطالبہ کیا۔ حضرت جعفرؓ نے نجاشی کے سامنے اسلام کا تعارف پیش کیا اور سورۃ
مریم کی ابتدائی آیات (حضرت عیسیٰ\,علیہ السلام کی ولادت کا بیان) تلاوت کیں۔
نجاشی رو پڑا اور بولا: ''یہ اور جو کچھ عیسیٰ لائے، ایک ہی چراغ سے نکلا ہے''
--- اور مہاجرین کو واپس کرنے سے انکار کر دیا۔

\medskip
\textbf{اہمیت۔} یہ اسلامی تاریخ کی پہلی ہجرت تھی، جس نے ثابت کیا کہ مسلمان
اپنے دین کی حفاظت کے لیے وطن چھوڑنے کو بھی تیار ہیں --- اور یہ مدینہ کی
بڑی ہجرت کی تمہید بنی۔
\end{hal}
